% =====================================================================
% City Planner stakeholder card (Chapter 6), area-data Geo design, N=4,000, seed 42.
% High case: Kingston upon Hull 029E (E01033104); low case: 003C (E01012877).
% Values: outputs/geo_area_data_eval/geo_coefficients_N4000.csv (signed beta, exact_w1),
%         outputs/geo_area_data_eval/spatial_patterns/ (Moran / LISA on signed beta),
%         outputs/feature_local_coefficients.csv (local Feature importance shares).
% "30x" is the ratio of TOTAL LOCAL FEATURE IMPORTANCE (152.4 vs 4.9), not of
% GeoSensitivity; the direct own-data measure is the exact single-LSOA shift (177 vs 1.1).
%
% Preamble requirements:
%   \usepackage[most]{tcolorbox}   \usepackage{xcolor}   \usepackage{graphicx}
% Copy into Chapter6/figs/: geo_area_data_map_N4000_plain.png (outputs/geo_area_data_eval/),
%   E01033104_feature_importance.png, E01012877_feature_importance.png
% =====================================================================

\definecolor{cardInk}{HTML}{12171C}
\definecolor{cardMuted}{HTML}{57646D}
\definecolor{cardTeal}{HTML}{1F4E5F}
\definecolor{cardPink}{HTML}{FCE4EC}
\definecolor{cardWine}{HTML}{8A0033}
\definecolor{cardRule}{HTML}{D9DEE2}

% Highlighted value, e.g. \pv{22.33}
\newcommand{\pv}[1]{\begingroup\setlength{\fboxsep}{1.5pt}\colorbox{cardPink}{\textcolor{cardWine}{\ttfamily\bfseries #1}}\endgroup}
% Role tag
\newcommand{\roletag}[1]{\tcbox[on line, colback=cardTeal, colframe=cardTeal, arc=2.5mm,
  boxsep=0pt, left=6pt, right=6pt, top=3pt, bottom=3pt]{\color{white}\sffamily\bfseries\small #1}}
% Section kicker, e.g. 01 · WHERE
\newcommand{\kicker}[1]{{\color{cardMuted}\sffamily\bfseries\footnotesize\MakeUppercase{#1}}}

\begin{figure}[htbp]
\begin{tcolorbox}[enhanced, colback=white, colframe=cardRule, arc=3mm, boxrule=0.6pt,
                  left=5mm, right=5mm, top=4mm, bottom=4mm, fontupper=\sffamily\small\color{cardInk}]

\roletag{CITY PLANNER}\par\medskip
{\large\itshape Where is the model most and least sensitive to local data?}\par\medskip

GeoSensitivity across Hull's \pv{168 LSOAs} is spatially clustered
(Moran's $I$ \pv{0.22}, $p = 0.001$).

\medskip
\begin{minipage}[c]{0.52\linewidth}
  \includegraphics[width=\linewidth]{Chapter6/figs/geo_area_data_map_N4000_plain.png}
\end{minipage}\hfill
\begin{minipage}[c]{0.44\linewidth}
  \kicker{01 \textperiodcentered{} Where}\par\smallskip
  One High--High cluster of \pv{9 LSOAs} (\pv{5.7\%} of Hull's area) runs from the
  Old Town along Hessle, Anlaby and Beverley Roads. Low--Low clusters sit on the edges,
  including Orchard Park.
\end{minipage}

\bigskip
\begin{minipage}[t]{0.485\linewidth}
  \kicker{02 \textperiodcentered{} High \textperiodcentered{} 029E, city centre}\par\smallskip
  \includegraphics[width=\linewidth]{Chapter6/figs/E01033104_feature_importance.png}\par\smallskip
  Rank \pv{1/168}, GeoSensitivity \pv{22.33}. Replacing its data with typical values
  shifts predictions by about \pv{177} crimes/month. Recent crime drives \pv{83\%} of
  local importance (6-month \pv{48.5}, 3-month \pv{37.8}).
\end{minipage}\hfill
\begin{minipage}[t]{0.485\linewidth}
  \kicker{03 \textperiodcentered{} Low \textperiodcentered{} 003C, Orchard Park}\par\smallskip
  \includegraphics[width=\linewidth]{Chapter6/figs/E01012877_feature_importance.png}\par\smallskip
  Rank \pv{146/168}, GeoSensitivity \pv{$-1.09$}. The same replacement shifts
  predictions by only about \pv{1.1}. Importance is spread: unemployment \pv{16\%},
  transit \pv{11\%}, crime history \pv{7\%}.
\end{minipage}

\bigskip
Total local Feature importance is \pv{30$\times$} larger in the city centre
(\pv{152.4} vs \pv{4.9}).

\medskip
{\color{cardRule}\rule{\linewidth}{0.4pt}}\par\smallskip
{\color{cardMuted}\footnotesize Higher GeoSensitivity = larger relative contribution to the model's
response; its sign is not a direction of predicted crime. Model behaviour, not planning targets or
causes of crime.}

\end{tcolorbox}
\caption{City Planner narrative: the most and a low GeoSensitivity area in Hull, with the
local Feature importance behind each.}
\label{fig:stakeholder-planner}
\end{figure}
